\documentclass[12pt]{article}
\usepackage{amsmath, amssymb, amsthm, geometry, mathtools}
\usepackage{hyperref}

\geometry{a4paper, margin=1in}

\title{Nonlinearity in Finite Dimensions as Linearity in Infinite-Dimensional Bundles: Koopman, Carleman, and Naimark Perspectives}
\author{Shimon Panfil}
\date{\today}

\theoremstyle{definition}
\newtheorem{definition}{Definition}[section]
\newtheorem{example}{Example}[section]

\theoremstyle{plain}
\newtheorem{theorem}{Theorem}[section]
\newtheorem{proposition}{Proposition}[section]

\begin{document}

\maketitle

\begin{abstract}
We unify several mathematical mechanisms by which nonlinear dynamics in finite-dimensional systems can be represented as linear evolution in infinite-dimensional spaces. Specifically, we formalize Koopman linearization, Carleman embedding, and Naimark dilation in the language of bundles and sections. This geometric viewpoint clarifies the deep structural analogy between classical flows, quantum measurements, and statistical evolution, highlighting the principle that nonlinearity is often an artifact of a reduced state description.
\end{abstract}

\section{Introduction}

Nonlinear dynamics in finite-dimensional spaces are ubiquitous in classical and quantum physics. However, it is often fruitful to reformulate such dynamics as \emph{linear} evolution in infinite-dimensional spaces. This idea underlies:

\begin{itemize}
    \item Koopman operator theory \cite{Koopman1931}, which linearizes classical flows via pullback on function spaces.
    \item Carleman embedding, which lifts polynomial ODEs to infinite-dimensional linear systems.
    \item Naimark dilation \cite{Naimark1940}, which represents POVMs as projective measurements in a larger Hilbert space.
\end{itemize}

We present a unified framework using the language of vector bundles and sections.

\section{Bundles and Sections}

\begin{definition}[Vector bundle]
A \emph{vector bundle} $\pi: E \to M$ over a smooth manifold $M$ is a smooth manifold $E$ such that each fiber $E_x := \pi^{-1}(x)$ is a vector space, and locally $E|_U \simeq U \times \mathbb{C}^k$.
\end{definition}

\begin{definition}[Section]
A \emph{section} of $E$ is a smooth map $s: M \to E$ with $\pi \circ s = \mathrm{id}_M$. The space of all smooth sections is denoted $\Gamma(E)$.
\end{definition}

Classical observables $f \in C^\infty(M)$ are sections of the trivial line bundle $M \times \mathbb{C}$, while quantum states can be modeled as sections of Hilbert bundles $\mathcal{H}_M \to M$.

\section{Koopman Linearization in Bundles}

Let $\Phi_t: M \to M$ be the flow generated by a smooth vector field $F \in \Gamma(TM)$. Define the pullback on sections:

\begin{equation}
U_t: \Gamma(E) \to \Gamma(E), \qquad U_t s := s \circ \Phi_t.
\end{equation}

\begin{theorem}[Bundle Koopman Linearization]
$U_t$ is linear and satisfies
\begin{equation}
U_{t+s} = U_t \circ U_s.
\end{equation}
Its generator is the Lie derivative along $F$:
\begin{equation}
L s = \mathcal{L}_F s.
\end{equation}
\end{theorem}

\begin{proof}
Linearity follows from the definition:
\[
U_t(a s_1 + b s_2)(x) = a s_1(\Phi_t(x)) + b s_2(\Phi_t(x)) = a U_t s_1(x) + b U_t s_2(x).
\]
The semigroup property is immediate from $\Phi_{t+s} = \Phi_t \circ \Phi_s$. The generator formula is standard from differentiation along the flow.
\end{proof}

\section{Carleman Linearization as Section Evolution}

Consider a polynomial ODE $\dot{x} = P(x)$, $x \in \mathbb{R}^n$. Define the symmetric tensor bundle:

\begin{equation}
S^\bullet(T^* \mathbb{R}^n) \to \mathbb{R}^n
\end{equation}

with sections $s = \sum_{k=0}^\infty s^{(k)}$ corresponding to monomials $x^{\otimes k}$. Then the ODE lifts to a linear operator on $\Gamma(S^\bullet)$:

\begin{equation}
\frac{d}{dt} s = \mathcal{L}_P s
\end{equation}

where $\mathcal{L}_P$ acts as a graded shift. Projection to degree-one monomials recovers the original nonlinear dynamics.

\section{Naimark Dilation and POVMs}

Let $\mathcal{H}$ be a Hilbert space and $E: \Sigma \to \mathcal{B}(\mathcal{H})$ a POVM. Naimark's theorem states that there exists a larger Hilbert space $\mathcal{K} \supset \mathcal{H}$, a PVM $\tilde{E}$ on $\mathcal{K}$, and an isometry $V: \mathcal{H} \to \mathcal{K}$ such that

\begin{equation}
E(A) = V^\dagger \tilde{E}(A) V, \quad \forall A \in \Sigma.
\end{equation}

This is a linearization of the nonlinear state-update associated with POVMs. Viewing $\mathcal{H}$ as sections of a Hilbert bundle, the PVM acts linearly on the extended bundle, and projection via $V^\dagger (\cdot) V$ recovers the original nonlinear collapse.

\section{Unified Perspective}

All three mechanisms share the same structural pattern:

\begin{enumerate}
    \item Start with nonlinear finite-dimensional dynamics or measurement.
    \item Embed into an infinite-dimensional bundle space.
    \item Evolution becomes linear on sections of this bundle.
    \item Projection back to original space reintroduces nonlinearity.
\end{enumerate}

Symbolically:
\[
\text{Nonlinear finite-dimensional dynamics} = \text{Projection of linear evolution in infinite-dimensional bundle}.
\]

\section{Discussion}

This viewpoint clarifies that nonlinearity is often a result of reduced state description. Linearization via bundles provides a unified geometric framework for:

\begin{itemize}
    \item Classical flows (Koopman),
    \item Polynomial ODEs (Carleman),
    \item Quantum measurements (Naimark).
\end{itemize}

It also suggests a deeper principle: the linear structure exists in a sufficiently large space; nonlinearity is induced by projecting onto a smaller subspace or fiber.

\section{Conclusion}

We presented a bundle-theoretic synthesis of linearization methods for nonlinear dynamics and measurements. Viewing dynamics as linear on sections of appropriately chosen bundles unifies classical and quantum perspectives, and highlights the geometric origin of linearity vs nonlinearity.

\begin{thebibliography}{9}

\bibitem{Koopman1931}
B. O. Koopman, \emph{Hamiltonian Systems and Transformation in Hilbert Space}, Proc. Natl. Acad. Sci. USA, 17 (1931), 315--318.

\bibitem{Naimark1940}
M. A. Naimark, \emph{Spectral Functions of a Symmetric Operator}, Izv. Akad. Nauk SSSR Ser. Mat. 4 (1940), 277–318.

\end{thebibliography}

\end{document}
